% Generated draft: requires mathematical and source review.
\documentclass[11pt]{article}
\usepackage[margin=25mm]{geometry}
\usepackage{fontspec}
\setmainfont{Latin Modern Roman}
\usepackage{amsmath,amssymb,mathtools,unicode-math}
\setmathfont{Latin Modern Math}
\usepackage{enumitem,needspace,xurl,hyperref}
\setlength{\parindent}{0pt}
\setlength{\parskip}{5pt}
\setlength{\emergencystretch}{3em}
\newcommand{\N}{\mathbb{N}}
\newcommand{\Z}{\mathbb{Z}}
\newcommand{\Q}{\mathbb{Q}}
\newcommand{\R}{\mathbb{R}}
\newcommand{\C}{\mathbb{C}}
\newcommand{\F}{\mathbb{F}}
\newcommand{\A}{\mathbb{A}}
\newcommand{\PP}{\mathbb{P}}
\newcommand{\E}{\mathbb{E}}
\newcommand{\eps}{\varepsilon}
\newcommand{\dd}{\,\mathrm{d}}
\begin{document}
\section*{UnsolvedMath ID: 800864}
\textbf{Source problem code:} catalogue entry
\subsection{Definitions and mathematical statement}
The supplied source asks for a geometric characterization of projective varieties of maximum likelihood degree one, but does not specify the ground field, irreducibility assumptions, or the convention for maximum likelihood degree. The following gives a precise formulation under the usual complex, embedded, irreducible-variety convention; these hypotheses require confirmation against the source context.

Let \(n\geq 0\), and let \(X\subseteq\PP^n_{\C}\) be a reduced irreducible projective variety with homogeneous coordinates \([p_0:\cdots:p_n]\). Set
\[
 s=p_0+\cdots+p_n,\qquad
 H=\{p_0p_1\cdots p_n s=0\},
\]
and assume \(X\not\subseteq H\). Write \(X_{\mathrm{reg}}\) for the smooth locus of \(X\), and put \(U=X_{\mathrm{reg}}\setminus H\). For a data vector \(u=(u_0,\ldots,u_n)\in\C^{n+1}\), define the logarithmic likelihood differential on \(U\) by
\[
 \omega_u=\sum_{i=0}^{n}u_i\,\frac{\mathrm d(p_i/s)}{p_i/s}.
\]
The ratios \(p_i/s\) are well-defined nonzero functions on \(U\). A likelihood critical point is a point \(x\in U\) such that \(\omega_u(x)\) vanishes on \(T_xU\). When the critical locus is finite for Zariski-generic \(u\), the maximum likelihood degree \(\operatorname{MLdeg}(X)\) is its generic number of points, counted with algebraic multiplicity.

\textbf{Problem (Sturmfels, 2007, Specific Problem 3).} Under this convention, find necessary and sufficient geometric conditions on the embedded variety \(X\subseteq\PP^n_{\C}\), with its specified coordinate hyperplanes and likelihood hyperplane \(s=0\), for
\[
 \operatorname{MLdeg}(X)=1.
\]
The source does not specify whether the intended characterization concerns this embedding and hyperplane configuration, the corresponding open variety, or an intrinsic property of the projective variety. It also does not specify how reducible varieties or varieties contained in \(H\) should be treated. Those extensions are unresolved in this draft.
\subsection{Short English statement}
Sturmfels asks for a geometric characterization of projective varieties whose maximum likelihood degree is one. Under the explicitly stated complex embedded-variety convention, this means that generic likelihood data have exactly one critical point, counted with multiplicity, on the smooth locus away from the coordinate and likelihood hyperplanes. The intended hypotheses and meaning of a geometric characterization require source-context review.
\subsection{Sources}
\begin{itemize}
\item Local UnsolvedMath registry: permanent ID \(800864\), exact source problem code \texttt{catalogue entry}, title \emph{Open Problems in Algebraic Statistics --- Specific Problem 3}. The registry attributes the problem to Sturmfels, proposed in 2007, and supplies the status \texttt{open}; that status is a supplied label, not an independently verified literature assessment.
\item Supplied primary-source reference: Sturmfels, \emph{Open Problems in Algebraic Statistics} (2007), Specific Problem \(\#3\), \url{https://arxiv.org/abs/0707.4558}. The supplied provenance records access on 2026-10-16 and extraction method \texttt{arxiv\_latex}; neither access nor the source text has been independently verified for this draft.
\item The supplied excerpt refers to \texttt{[CHKS, HKS]}, but provides no bibliographic entries identifying these references.
\item Supplied upstream revision: \texttt{d371ee917ac9647fbb06cd6fc529deec6f3dda82}.
\end{itemize}
\end{document}
